% !TeX root = ../../Book3.tex
\section{Laurent 级数}
\label{b3:sec:0301}

Taylor 级数要求中心点也在定义域内。去掉中心后，负幂可能出现，但它们并不破坏环域内部的一致收敛。把内外两条边界上的 Cauchy 核分别展开，就能同时得到正幂与负幂；两个部分有不同的几何来源。

\subsection{环域}
\label{b3:subsec:030101}

设 \(0\le r<R\le\infty\)，记
\[
 A(a;r,R)\coloneqq\{\,z\in\mathbb C\mid r<|z-a|<R\,\}.
\]
它称为以 \(a\) 为中心的\term[huanyu]{环域}[annulus]；\(r=0\) 时是去心圆盘。所谓 Laurent 级数的收敛，指正幂部分与负幂部分分别收敛，不能用正负指标混合的任意排列改变和。

\ProseParagraphBreak

负幂级数 \(\sum_{n\ge1}c_{-n}(z-a)^{-n}\) 在倒数变量 \(w=1/(z-a)\) 中是没有常数项的幂级数。若该幂级数半径为 \(S\)，它在 \(|z-a|>1/S\) 上收敛；正幂部分则在 \(|z-a|<R\) 上收敛。两者的公共范围因而天然是环域，可能为空。边界圆上的收敛须分别考察。

\ProseParagraphBreak

例如 \(1/(z(z-1))\) 在 \(0<|z|<1\) 时写成 \(-\sum_{n\ge0}z^{n-1}\)，在 \(|z|>1\) 时写成 \(\sum_{n\ge0}z^{-n-2}\)。展开中心相同，而穿过 \(|z|=1\) 时表达式改变，原因是圆周上的奇点 \(z=1\) 阻断了统一展开。

\ProseParagraphBreak

选择展开环域时，可以先将所有与中心的距离画在数轴上。例如函数 \(1/((z-1)(z-2))\) 的奇点到中心0的距离为1与2，因此有三个候选范围：\(|z|<1\)、\(1<|z|<2\) 和 \(|z|>2\)。中间范围中的 \(1/(z-1)\) 应展开成负幂，而 \(1/(z-2)\) 应展开成正幂。同一个有理式里既有正幂又有负幂，并不表示使用了互相矛盾的几何级数；它们的公共收敛域正好是指定环域。

\ProseParagraphBreak

Laurent 主部是否有限，也取决于中心在哪里。函数 \(1/(z-1)\) 在中心0的小圆盘是 Taylor 级数，在中心1则是只有一个负幂的 Laurent 级数，在无穷远坐标又具有可去奇点。分类和系数总要附一个中心，不能只说“这个函数有负幂”而略去展开范围。

\ProseParagraphBreak

若两个奇点到中心距离相同，例如 \(1/(z^2+1)\)，它们在同一分界圆上。尽管圆周上只有两个点坏掉，也不能把圆内与圆外的 Laurent 展开拼成一个穿过该圆的同心环域展开，因为一个 Laurent 级数在其开放收敛环域上处处全纯，会把这两个真极点也包括进去。

\subsection{Laurent 展开}
\label{b3:subsec:030102}

\term[laurent-jishu]{Laurent 级数}[Laurent series]允许整数幂指标。
\begin{theorem}[Laurent]\label{b3:thm:03-laurent}
若 \(f\) 在 \(A(a;r,R)\) 上全纯，则有唯一展开
\[
 f(z)=\sum_{n=-\infty}^{\infty}c_n(z-a)^n,
\]
其中正幂部分在 \(|z-a|<R\) 全纯，负幂部分在 \(|z-a|>r\) 全纯；级数在每个闭子环域上一致绝对收敛。
\end{theorem}
\begin{proof}
固定 \(r<\rho_1<|z-a|<\rho_2<R\)。沿环域的一条径向线切开，将内外圆的边界用两条相反方向的切口边连接；切口可以选择不经过 \(z\)。对切开区域使用 Cauchy 公式，再让两条切口边靠拢。核与函数在切口附近连续，两个切口积分相消，得到
\[
 f(z)=\frac1{2\pi\mathrm{i}}\int_{|\zeta-a|=\rho_2}\frac{f(\zeta)}{\zeta-z}\,\mathrm{d}\zeta-
 \frac1{2\pi\mathrm{i}}\int_{|\zeta-a|=\rho_1}\frac{f(\zeta)}{\zeta-z}\,\mathrm{d}\zeta,
\]
两圆均取逆时针方向。在外圆上，展开 \((\zeta-z)^{-1}=\sum_{n\ge0}(z-a)^n/(\zeta-a)^{n+1}\)；在内圆上，展开 \((\zeta-z)^{-1}=-\sum_{n\ge0}(\zeta-a)^n/(z-a)^{n+1}\)。若 \(z\) 限于严格夹在两圆之间的闭子环域，两个几何级数都一致绝对收敛，故可以逐项积分。

所得系数为 \(c_n=(2\pi\mathrm{i})^{-1}\int_{|\zeta-a|=\rho}f(\zeta)/(\zeta-a)^{n+1}\,\mathrm{d}\zeta\)。在两个容许半径之间，被积函数全纯；对其环形边界切开并用 Cauchy 定理，两个圆周积分相等。因而系数不依赖所选半径。取半径分别靠近 \(r\) 和 \(R\)，前述表示覆盖整个环域。对任意闭子环域，选更内、更外两圆，几何估计给出一致绝对收敛。

若另有一组 Laurent 系数，在某个中间圆周上逐项积分乘以 \((z-a)^{-k-1}\)，只有指标 \(k\) 的项留下 \(2\pi\mathrm{i}\)。于是两组系数逐个相等。
\end{proof}
这里切口不是对函数选分支，而是把有孔区域分解成可以使用 Cauchy 公式的区域。切口两侧的函数本来相同，故会相消；对多值表达式不能直接照用这一步。

\subsection{Laurent 系数}
\label{b3:subsec:030103}

\begin{proposition}[Laurent 系数公式]\label{b3:prop:03-laurent-coefficients}
在 Laurent 定理的条件下，对任意 \(r<\rho<R\) 与整数 \(n\)，
\[
 c_n=\frac1{2\pi\mathrm{i}}\int_{|z-a|=\rho}\frac{f(z)}{(z-a)^{n+1}}\,\mathrm{d}z,
 \qquad |c_n|\le M_\rho\rho^{-n},
\]
其中 \(M_\rho=\max_{|z-a|=\rho}|f(z)|\)。
\end{proposition}
\begin{proof}
公式已由前一定理的构造与唯一性得到。对积分使用长度估计，圆长 \(2\pi\rho\) 与分母 \(\rho^{n+1}\) 相消，即得对正负指标同时有效的估计。
\end{proof}
当 \(n=-m<0\) 时，估计成为 \(|c_{-m}|\le M_\rho\rho^m\)。若函数靠近中心有界，令 \(\rho\downarrow0\) 就迫使全部负幂系数为零。正幂系数由外侧增长控制，负幂系数由内侧增长控制；同一公式的两个方向将给出奇点分类。

\ProseParagraphBreak

级数仍可在闭子环域上逐项微分，因为两部分分别是在原变量与倒数变量中的幂级数；微分后 \((z-a)^n\) 产生 \(n(z-a)^{n-1}\)，包括负指标。逐项原函数则需额外检查 \(n=-1\) 项，它不能在整个环域上取单值原函数。

\ProseParagraphBreak

Laurent 展开中的系数可由对称性提前筛选。若 \(f(-z)=f(z)\) 且定义环域以0为中心，唯一性给 \(c_n=(-1)^n c_n\)，所以奇数指标系数为零，包括负的奇数指标。若 \(f(\overline z)=\overline{f(z)}\)，则把展开共轭并替换变量后，由唯一性得到全部系数为实数。这里的对称必须在整个环域成立；只在实轴知道函数值实，并且有内部聚点时，也可由恒等定理传播，但要先保证所比较的共轭反射函数全纯。

\ProseParagraphBreak

负幂项不能用普通原函数遗漏。若 \(f\) 在同心环域的留数为0，可以对全部其他项逐项积分，得到在该环域收敛的 Laurent 原函数。其正幂与负幂部分仍在闭子环域一致收敛，可用幂级数估计核对。若留数非零，绕中心圆周积分非零，端点公式立即排除原函数。因此对同心环域上的全纯函数，留数为零恰是存在单值原函数的条件；对于有多个孔洞的区域，要分别检查各孔的周期。

\begin{proposition}\label{b3:prop:03-annulus-primitive}
设 \(f\) 在环域 \(r<|z-a|<R\) 全纯，Laurent 系数为 \(c_n\)。它在该环域有单值全纯原函数，当且仅当 \(c_{-1}=0\)。
\end{proposition}
\begin{proof}
有原函数时任意中间圆周的积分为零，所以 \(2\pi\mathrm{i}c_{-1}=0\)。反之，去掉 \(n=-1\) 项后令
\[
 F(z)=\sum_{n\ne-1}\frac{c_n}{n+1}(z-a)^{n+1}.
\]
在任意闭子环域上，正幂与负幂部分均正规收敛；略扩大子环域后用几何估计可同时控制导数级数。因此逐项求导合法，且 \(F'=f\)。
\end{proof}
这里每个负幂项 \((z-a)^{-n}\)、\(n\ge2\) 都有原函数；只有 \((z-a)^{-1}\) 会产生对数的多值性。故“奇点存在”并不等于“没有原函数”。\(1/z^2\) 虽在0有极点，却在去心平面有原函数 \(-1/z\)。闭路积分测量的是留数这一特定系数，不能直接读出全部主部。
